\section{실험 결과}

\subsection*{실패 발견율과 Random 대비 향상}
\begin{table}[!t]
\centering\footnotesize\setlength{\tabcolsep}{3pt}
\caption{AFS와 Random의 실패 발견 결과}\label{tab:result}
\begin{tabularx}{\linewidth}{@{}Ycccc@{}}
\toprule
& \multicolumn{2}{c}{실험 1: 로봇팔} & \multicolumn{2}{c}{실험 2: 휴머노이드} \\
\cmidrule(lr){2-3}\cmidrule(lr){4-5}
& AFS & Random & AFS & Random \\
\midrule
유효 시도 & 12 & 12 & 5 & 5 \\
PASS / FAIL & 5 / 7 & 10 / 2 & 3 / 2 & 5 / 0 \\
\textbf{실패 발견율 (FDR)} & \textbf{58.3\%} & 16.7\% & \textbf{40.0\%} & 0.0\% \\
탐색 단계 FDR (초기 장면 제외) & 6/6 (100\%) & 1/6 (16.7\%) & 2/4 (50\%) & 0/4 (0\%) \\
실패한 고유 환경 수 & 4 & 2 & 2 & 0 \\
\bottomrule
\end{tabularx}
\end{table}

\begin{figure}[!t]
\centering
\begin{tikzpicture}
\begin{axis}[
  ybar, width=\linewidth, height=48mm, bar width=9pt,
  ymin=0, ymax=125, ylabel={FDR (\%)}, ylabel style={font=\small},
  symbolic x coords={A,B,C,D}, xtick=data,
  xticklabels={{로봇팔\\전체},{로봇팔\\탐색 단계},{휴머노이드\\전체},{휴머노이드\\탐색 단계}},
  xticklabel style={font=\scriptsize,align=center},
  tick label style={font=\scriptsize},
  nodes near coords, nodes near coords style={font=\scriptsize},
  legend style={font=\scriptsize,at={(0.5,1)},anchor=north,draw=none,legend columns=2},
  axis lines*=left, enlarge x limits=0.18]
\addplot[fill=afs,draw=afs] coordinates {(A,58.3) (B,100) (C,40) (D,50)};
\addplot[fill=rnd,draw=rnd] coordinates {(A,16.7) (B,16.7) (C,0) (D,0)};
\legend{AFS, Random}
\end{axis}
\end{tikzpicture}
\caption{방법별 실패 발견율. 탐색 단계는 공통 초기 장면을 제외한 적응적 선택 구간이다.}\label{fig:fdr}
\end{figure}

표~\ref{tab:result}와 그림~\ref{fig:fdr}에 결과를 정리하였다. 로봇팔 실험에서 AFS의 FDR은 58.3\%로 목표 30\%를 넘었고, Random(16.7\%) 대비 상대 향상률은
$(58.3-16.7)/16.7 = 250\%$로 목표 20\%를 크게 웃돈다. 공통 초기 장면을 제외한 탐색 단계만 보면
AFS는 6회 모두 실패를 찾았고 Random은 1회였다. 시드별로도 AFS는 3/2/2건, Random은 1/1/0건으로
모든 시드에서 AFS가 더 많은 실패를 발견했다. 휴머노이드 실험에서도 AFS는 40\%, Random은 0\%로
40퍼센트포인트 차이를 보였다. 사람 안전 이벤트는 AFS 4/5회에서 검출되었다.

\subsection*{실패 유형 다양성}
\begin{table*}[!t]
\centering\small
\caption{초기 정의 6종 기준 실패 유형 발견 현황}\label{tab:div}
\begin{tabularx}{\textwidth}{@{}p{24mm}cp{26mm}Y@{}}
\toprule
유형 & 발견 & 실험 & 관측 근거 \\
\midrule
충돌 & \checkmark & 로봇팔 & 로봇–장애물 접촉 10.7--13.7초, 최대 접촉력 약 200--280\,N, 그릇 이동 1\,mm 미만 \\
장애물 간섭 & \checkmark & 휴머노이드 & 장애물 부근에서 경로 계획이 반복 거절되며 정체 \\
목표 위치 점유 & \checkmark & 휴머노이드 & 상자가 목표 영역을 점유, 목표 0.86\,m 앞에서 정체 \\
작업자 안전 위험 & \checkmark & 휴머노이드 & 사람과 1\,m 이내 0.5초 이상 노출 또는 접촉, AFS 4/5회 \\
도달 불가 & 후보 & 로봇팔·휴머노이드 & 그릇을 들어 운반했으나 접시 8.7\,cm 앞에서 미완료, 최종 목표점 경로 계획 반복 거절 \\
인식 오류 & -- & -- & 가림·유사 물체 조건은 후속 실험 \\
\midrule
\multicolumn{4}{@{}l}{\textbf{합계: 6종 중 4종 발견 + 도달 불가 후보 1종 (목표 4종 달성)}} \\
\bottomrule
\end{tabularx}
\end{table*}
표~\ref{tab:div}는 초기 정의 6종 기준의 발견 현황이다. 작업자 안전 위험은 목표 달성 여부와 무관하게 안전 요구사항을 위반한 실행을 실패 사례로 본다.
실제 실행에서는 초기 6종으로 설명하기 어려운 행동 실패도 관측되어, 넘어짐·이동 정체·물체 조작
실패·목표 완료 실패와 시스템 연동 문제 2종(최종 목표점 경로 계획 거절, 중간 목표 반복 도착)을 더한 12종 체계로 확장하였다(부록 B).
이 확장 체계로는 목표 완료 실패, 물체 조작 실패, 중간 목표 반복을 추가로 관측하였다.
